% File: ai-nhap/2026-11-15_11_de_tham_khao_quan_11_de_2_va_quan_12_de_1.tex
% Nguồn: Tuyển tập đề tham khảo môn Toán Tuyển sinh 10 TP.HCM 2025-2026
% Gồm 14 bài: Quận 11 - Đề 2 (7 bài) và Quận 12 - Đề 1 (7 bài)

% ==============================================================================
% PHẦN 1: ĐỀ THAM KHẢO TUYỂN SINH 10 QUẬN 11 – ĐỀ 2 (7 BÀI)
% ==============================================================================

\begin{bt}%[9D6H18-2]
% nguon: Đề tham khảo Tuyển sinh 10 năm học 2025-2026 Sở GD&ĐT TP.HCM – Quận 11 - Đề 2 (Bài 1)
Cho parabol $(P)\colon y = -\dfrac{x^2}{4}$.
\begin{enumerate}[a)]
    \item Vẽ đồ thị $(P)$ trên hệ trục tọa độ $Oxy$.
    \item Tìm tọa độ các điểm $M$ thuộc đồ thị $(P)$ sao cho tung độ bằng hai lần hoành độ.
\end{enumerate}
\loigiai{
\begin{enumerate}[a)]
    \item Bảng giá trị tương ứng của hàm số $y = -\dfrac{x^2}{4}$:
    \begin{center}
    \begin{tabular}{|c|c|c|c|c|c|}
        \hline
        $x$ & $-4$ & $-2$ & $0$ & $2$ & $4$ \\
        \hline
        $y = -\dfrac{x^2}{4}$ & $-4$ & $-1$ & $0$ & $-1$ & $-4$ \\
        \hline
    \end{tabular}
    \end{center}
    Đồ thị $(P)$ là một parabol có đỉnh tại gốc tọa độ $O(0; 0)$, nhận trục tung $Oy$ làm trục đối xứng và quay bề lõm xuống phía dưới.
    \item Điểm $M(x; y)$ thuộc $(P)$ có tung độ bằng hai lần hoành độ khi và chỉ khi
    $y = 2x \iff -\dfrac{x^2}{4} = 2x \iff x^2 + 8x = 0 \iff x(x + 8) = 0 \iff \left[\begin{aligned} &x = 0 \implies y = 0 \\ &x = -8 \implies y = 2 \cdot (-8) = -16. \end{aligned}\right.$
    Vậy có hai điểm thỏa mãn yêu cầu đề bài là $M_1(0; 0)$ và $M_2(-8; -16)$.
\end{enumerate}
}
\end{bt}

\begin{bt}%[9D6H20-2]
% nguon: Đề tham khảo Tuyển sinh 10 năm học 2025-2026 Sở GD&ĐT TP.HCM – Quận 11 - Đề 2 (Bài 2)
Cho phương trình bậc hai $3x^2 - 6x - 7 = 0$.
\begin{enumerate}[a)]
    \item Chứng minh phương trình luôn có hai nghiệm phân biệt $x_1, x_2$.
    \item Không giải phương trình, hãy tính giá trị của biểu thức $A = \dfrac{6x_1^2 - 14}{x_2} + \dfrac{6x_2^2 - 14}{x_1}$.
\end{enumerate}
\loigiai{
\begin{enumerate}[a)]
    \item Phương trình $3x^2 - 6x - 7 = 0$ có các hệ số $a = 3$, $c = -7$.
    Vì $a \cdot c = 3 \cdot (-7) = -21 < 0$ nên phương trình luôn có hai nghiệm phân biệt trái dấu $x_1, x_2$.
    \item Vì $x_1, x_2$ là các nghiệm của phương trình $3x^2 - 6x - 7 = 0$ nên:
    $\begin{cases} 3x_1^2 - 6x_1 - 7 = 0 \implies 3x_1^2 - 7 = 6x_1 \implies 6x_1^2 - 14 = 12x_1 \\ 3x_2^2 - 6x_2 - 7 = 0 \implies 3x_2^2 - 7 = 6x_2 \implies 6x_2^2 - 14 = 12x_2. \end{cases}$
    Thay vào biểu thức $A$, ta được:
    $A = \dfrac{12x_1}{x_2} + \dfrac{12x_2}{x_1} = 12 \left( \dfrac{x_1}{x_2} + \dfrac{x_2}{x_1} \right) = 12 \left( \dfrac{x_1^2 + x_2^2}{x_1 x_2} \right) = 12 \cdot \dfrac{(x_1 + x_2)^2 - 2x_1 x_2}{x_1 x_2}$.
    Theo định lí Viète:
    $\begin{cases} S = x_1 + x_2 = -\dfrac{b}{a} = \dfrac{6}{3} = 2 \\ P = x_1 x_2 = \dfrac{c}{a} = -\dfrac{7}{3}. \end{cases}$
    Do đó:
    $(x_1 + x_2)^2 - 2x_1 x_2 = 2^2 - 2 \cdot \left(-\dfrac{7}{3}\right) = 4 + \dfrac{14}{3} = \dfrac{26}{3}$.
    Suy ra:
    $A = 12 \cdot \dfrac{\frac{26}{3}}{-\frac{7}{3}} = 12 \cdot \left(-\dfrac{26}{7}\right) = -\dfrac{312}{7}$.
\end{enumerate}
}
\end{bt}

\begin{bt}%[9D6V21-3]
% nguon: Đề tham khảo Tuyển sinh 10 năm học 2025-2026 Sở GD&ĐT TP.HCM – Quận 11 - Đề 2 (Bài 3)
Một sân chơi hình chữ nhật có chiều dài là $x$ mét và chiều dài hơn chiều rộng là $3\text{ m}$.
\begin{enumerate}[a)]
    \item Viết biểu thức tính diện tích $S$ của sân chơi hình chữ nhật theo $x$.
    \item Sau đó, ban quản lý quyết định tăng chiều dài thêm $1\text{ m}$ và tăng chiều rộng thêm $3\text{ m}$ thì diện tích của sân chơi hình chữ nhật tăng gấp đôi. Tính chiều dài ban đầu của sân chơi đó.
\end{enumerate}
\loigiai{
\begin{enumerate}[a)]
    \item Chiều rộng của sân chơi hình chữ nhật ban đầu là $x - 3$ (m), điều kiện $x > 3$.
    Diện tích ban đầu của sân chơi là:
    $S = x(x - 3) = x^2 - 3x\text{ (m}^2\text{)}$.
    \item Chiều dài mới sau khi tăng thêm $1\text{ m}$ là: $x + 1$ (m).
    Chiều rộng mới sau khi tăng thêm $3\text{ m}$ là: $(x - 3) + 3 = x$ (m).
    Diện tích mới của sân chơi là: $(x + 1)x = x^2 + x\text{ (m}^2\text{)}$.
    Vì diện tích mới tăng gấp đôi diện tích ban đầu nên ta có phương trình:
    $x^2 + x = 2(x^2 - 3x) \iff x^2 + x = 2x^2 - 6x \iff x^2 - 7x = 0 \iff x(x - 7) = 0$.
    Vì $x > 3$ nên $x = 7$ (thỏa mãn điều kiện).
    Vậy chiều dài ban đầu của sân chơi là $7\text{ m}$.
\end{enumerate}
}
\end{bt}

\begin{bt}%[9H0V31-3]
% nguon: Đề tham khảo Tuyển sinh 10 năm học 2025-2026 Sở GD&ĐT TP.HCM – Quận 11 - Đề 2 (Bài 4)
Để phòng tránh tai nạn đuối nước cho trẻ em, người ta quyết định dùng đất để san lấp một cái ao có dạng nửa hình cầu, mặt ao là hình tròn có đường kính $10\text{ m}$.
\begin{enumerate}[a)]
    \item Tính thể tích nước chứa trong ao theo $\text{m}^3$ khi mặt nước ngang bằng mặt đất (kết quả làm tròn đến hàng phần mười).
    \item Người ta thuê các xe tải có lòng thùng dạng hình hộp chữ nhật với chiều dài $9{,}9\text{ m}$, chiều rộng $2{,}37\text{ m}$ và chiều cao $0{,}85\text{ m}$. Do tuyến đường vận chuyển bị giới hạn tải trọng nên mỗi xe chỉ được chở tối đa $85\%$ thể tích lòng thùng xe. Hỏi cần thuê ít nhất bao nhiêu chuyến xe để vận chuyển đủ đất san lấp đầy cái ao? (Giả sử thể tích đất san lấp bằng thể tích lòng ao).
\end{enumerate}
\loigiai{
\begin{enumerate}[a)]
    \item Bán kính của ao nước hình bán cầu là $R = \dfrac{10}{2} = 5\text{ m}$.
    Thể tích nước trong ao bằng một nửa thể tích của khối cầu bán kính $R$:
    $V_{\text{ao}} = \dfrac{1}{2} \cdot \left(\dfrac{4}{3}\pi R^3\right) = \dfrac{2}{3}\pi \cdot 5^3 = \dfrac{250\pi}{3} \approx 261{,}8\text{ m}^3$.
    \item Thể tích toàn bộ lòng trong của thùng xe tải là:
    $V_{\text{thùng}} = 9{,}9 \cdot 2{,}37 \cdot 0{,}85 = 19{,}94355\text{ m}^3$.
    Mỗi chuyến xe chỉ được phép chở tối đa $85\%$ thể tích thùng:
    $V_{\text{chở}} = 0{,}85 \cdot V_{\text{thùng}} = 0{,}85 \cdot 19{,}94355 = 16{,}9520175\text{ m}^3$.
    Số chuyến xe cần ít nhất để chở đủ lượng đất san lấp ao là:
    $n = \dfrac{V_{\text{ao}}}{V_{\text{chở}}} = \dfrac{261{,}8}{16{,}9520175} \approx 15{,}44$.
    Vì số chuyến xe phải là số nguyên dương nên cần thuê ít nhất $16$ chuyến xe tải.
\end{enumerate}
}
\end{bt}

\begin{bt}%[9D1H3-4]
% nguon: Đề tham khảo Tuyển sinh 10 năm học 2025-2026 Sở GD&ĐT TP.HCM – Quận 11 - Đề 2 (Bài 5)
Theo nghiên cứu dinh dưỡng, lượng protein chứa trong thịt bò và thịt gà lần lượt là $26\%$ và $31\%$; lượng lipit chứa trong thịt bò và thịt gà lần lượt là $15\%$ và $14\%$. Lượng protein và lipit tiêu chuẩn cần thiết cho người trưởng thành mỗi ngày lần lượt là $0{,}8\text{ g/kg}$ và $0{,}9\text{ g/kg}$ trọng lượng cơ thể. Bác Dũng có trọng lượng $70\text{ kg}$. Hỏi bác Dũng cần bổ sung bao nhiêu gam protein và lipit mỗi ngày theo trọng lượng cơ thể? Nếu chỉ dùng thịt bò để bổ sung đủ lượng protein thì cần bao nhiêu gam thịt bò (làm tròn đến hàng đơn vị)?
\loigiai{
Lượng protein cần thiết cho bác Dũng mỗi ngày là:
$m_{\text{protein}} = 70 \cdot 0{,}8 = 56\text{ g}$.
Lượng lipit cần thiết cho bác Dũng mỗi ngày là:
$m_{\text{lipit}} = 70 \cdot 0{,}9 = 63\text{ g}$.
Nếu bác Dũng chỉ ăn thịt bò để bổ sung đủ $56\text{ g}$ protein thì lượng thịt bò cần dùng là:
$m_{\text{bò}} = \dfrac{56}{26\%} = \dfrac{56}{0{,}26} \approx 215\text{ g}$.
}
\end{bt}

\begin{bt}%[9D8H26-3]
% nguon: Đề tham khảo Tuyển sinh 10 năm học 2025-2026 Sở GD&ĐT TP.HCM – Quận 11 - Đề 2 (Bài 6)
Bạn An tham gia một trò chơi ném bi tại hội chợ xuân. Đích ném là một bảng gồm $25$ ô vuông bằng nhau, điểm số của mỗi ô được quy định như bảng sau:
\begin{center}
\begin{tabular}{|c|c|c|c|c|}
    \hline
    $5$ & $3$ & $3$ & $3$ & $5$ \\
    \hline
    $3$ & $-2$ & $-1$ & $-2$ & $3$ \\
    \hline
    $3$ & $-1$ & $5$ & $-1$ & $3$ \\
    \hline
    $3$ & $-2$ & $-1$ & $-2$ & $3$ \\
    \hline
    $5$ & $3$ & $3$ & $3$ & $5$ \\
    \hline
\end{tabular}
\end{center}
Luật chơi: Mỗi lần ném trúng vào một ô thì nhận được số điểm ghi trong ô đó, nếu ném trượt ra ngoài bảng thì bị trừ $5$ điểm. Sau $10$ lần ném: đạt từ $30$ điểm đến $49$ điểm nhận được phần quà trị giá $300\,000$ đồng.
\begin{enumerate}[a)]
    \item Trong $9$ lần ném đầu tiên, bạn An ném được $5$ lần vào ô điểm $5$, $1$ lần trượt ra ngoài bảng, $2$ lần vào ô điểm $3$, $1$ lần vào ô điểm $-1$. Tính tổng số điểm bạn An đạt được sau $9$ lần ném.
    \item Hỏi bạn An có cơ hội nhận phần quà trị giá $300\,000$ đồng hay không? Nếu có thì ở lần ném thứ $10$ bạn An cần ném vào ô điểm nào? Tính xác suất để bạn An nhận được phần quà đó (giả sử khả năng ném vào mỗi ô trong $25$ ô là như nhau).
\end{enumerate}
\loigiai{
\begin{enumerate}[a)]
    \item Sau $9$ lần ném, tổng số điểm của bạn An là:
    $S_9 = 5 \cdot 5 + 1 \cdot (-5) + 2 \cdot 3 + 1 \cdot (-1) = 25 - 5 + 6 - 1 = 25$ (điểm).
    \item Để nhận được phần quà trị giá $300\,000$ đồng thì sau $10$ lần ném tổng điểm phải từ $30$ điểm đến $49$ điểm.
    Số điểm bạn An cần đạt thêm ở lần ném thứ $10$ là:
    $x \ge 30 - 25 = 5$ (và $x \le 49 - 25 = 24$).
    Quan sát bảng điểm gồm $25$ ô, điểm số cao nhất của mỗi ô là $5$. Do đó để đạt được tổng điểm ít nhất $30$ điểm, ở lần ném thứ $10$ bạn An bắt buộc phải ném trúng vào ô điểm $5$.
    Trên bảng có tất cả $25$ ô, trong đó có $5$ ô mang điểm $5$ (gồm $4$ ô ở $4$ góc và $1$ ô ở chính giữa).
    Do đó bạn An hoàn toàn có cơ hội nhận phần quà trị giá $300\,000$ đồng khi ném vào ô điểm $5$.
    Xác suất để bạn An nhận được phần quà đó là:
    $P = \dfrac{5}{25} = \dfrac{1}{5} = 0{,}2$.
\end{enumerate}
}
\end{bt}

\begin{bt}%[9H9C29-3]
% nguon: Đề tham khảo Tuyển sinh 10 năm học 2025-2026 Sở GD&ĐT TP.HCM – Quận 11 - Đề 2 (Bài 7)
Cho đường tròn $(O)$ và điểm $A$ nằm bên ngoài đường tròn. Vẽ hai tiếp tuyến $AM, AN$ với đường tròn $(O)$ ($M, N$ là các tiếp điểm). Gọi $H$ là giao điểm của dây $MN$ và đoạn thẳng $OA$.
\begin{enumerate}[a)]
    \item Chứng minh tứ giác $AMON$ nội tiếp đường tròn.
    \item Kẻ dây cung $BC$ của đường tròn $(O)$ đi qua điểm $H$ sao cho $BC \parallel AM$ ($C$ thuộc cung nhỏ $MN$). Đường thẳng $AC$ cắt đường tròn $(O)$ tại điểm thứ hai là $D$. Gọi $I$ là trung điểm của dây $CD$ và $F$ là giao điểm của $MN$ và $CD$. Chứng minh $\widehat{NMI} = \widehat{NOI}$.
    \item Chứng minh $FI \cdot FA = FC \cdot FD$.
\end{enumerate}
\loigiai{
\begin{enumerate}[a)]
    \item Vì $AM, AN$ là hai tiếp tuyến của $(O)$ nên $AM \perp OM$ và $AN \perp ON$, suy ra $\widehat{AMO} = \widehat{ANO} = 90^\circ$.
    Tứ giác $AMON$ có $\widehat{AMO} + \widehat{ANO} = 90^\circ + 90^\circ = 180^\circ$.
    Do đó tứ giác $AMON$ nội tiếp đường tròn đường kính $AO$.
    \item Vì $I$ là trung điểm của dây cung $CD$ không đi qua tâm $O$ nên $OI \perp CD$ tại $I$, suy ra $\widehat{AIO} = 90^\circ$.
    Từ đó điểm $I$ thuộc đường tròn đường kính $AO$, nên năm điểm $A, M, I, O, N$ cùng thuộc một đường tròn.
    Trong đường tròn đường kính $AO$ ngoại tiếp ngũ giác $AMION$, hai góc $\widehat{NMI}$ và $\widehat{NOI}$ là hai góc nội tiếp cùng chắn cung $NI$.
    Do đó $\widehat{NMI} = \widehat{NOI}$ (đpcm).
    \item Vì bốn điểm $M, I, O, N$ cùng thuộc đường tròn ngoại tiếp tứ giác $AMON$ và $F$ là giao điểm của hai dây cung $MN$ và $CD$ kéo dài, ta có $\triangle FMI \backsim \triangle FON$ (g.g), suy ra $FM \cdot FN = FI \cdot FO$ hoặc áp dụng phương tích:
    Trong đường tròn $(O)$, hai cát tuyến $FMN$ và $FCD$ cắt nhau tại $F$ nên $FM \cdot FN = FC \cdot FD$.
    Mặt khác, xét đường tròn đường kính $AO$ đi qua bốn điểm $A, I, H, O$, ta có hệ thức $FI \cdot FA = FM \cdot FN$.
    Từ đó suy ra $FI \cdot FA = FC \cdot FD$ (đpcm).
\end{enumerate}
}
\end{bt}


% ==============================================================================
% PHẦN 2: ĐỀ THAM KHẢO TUYỂN SINH 10 QUẬN 12 – ĐỀ 1 (7 BÀI)
% ==============================================================================

\begin{bt}%[9D6H18-2]
% nguon: Đề tham khảo Tuyển sinh 10 năm học 2025-2026 Sở GD&ĐT TP.HCM – Quận 12 - Đề 1 (Bài 1)
Cho parabol $(P)\colon y = -\dfrac{1}{2}x^2$.
\begin{enumerate}[a)]
    \item Vẽ đồ thị $(P)$ trên hệ trục tọa độ $Oxy$.
    \item Tìm tọa độ các điểm thuộc $(P)$ có hoành độ gấp hai lần tung độ.
\end{enumerate}
\loigiai{
\begin{enumerate}[a)]
    \item Bảng giá trị tương ứng của hàm số $y = -\dfrac{1}{2}x^2$:
    \begin{center}
    \begin{tabular}{|c|c|c|c|c|c|}
        \hline
        $x$ & $-4$ & $-2$ & $0$ & $2$ & $4$ \\
        \hline
        $y = -\dfrac{1}{2}x^2$ & $-8$ & $-2$ & $0$ & $-2$ & $-8$ \\
        \hline
    \end{tabular}
    \end{center}
    Đồ thị $(P)$ là một parabol có đỉnh tại gốc tọa độ $O(0; 0)$, nhận trục $Oy$ làm trục đối xứng và quay bề lõm xuống phía dưới.
    \item Điểm $M(x; y)$ thuộc $(P)$ có hoành độ gấp hai lần tung độ khi và chỉ khi:
    $x = 2y \iff y = \dfrac{x}{2}$.
    Thay vào phương trình parabol $(P)$, ta có:
    $-\dfrac{1}{2}x^2 = \dfrac{x}{2} \iff x^2 + x = 0 \iff x(x + 1) = 0 \iff \left[\begin{aligned} &x = 0 \implies y = 0 \\ &x = -1 \implies y = -\dfrac{1}{2}. \end{aligned}\right.$
    Vậy có hai điểm thỏa mãn yêu cầu là $(0; 0)$ và $\left(-1; -\dfrac{1}{2}\right)$.
\end{enumerate}
}
\end{bt}

\begin{bt}%[9D6H20-2]
% nguon: Đề tham khảo Tuyển sinh 10 năm học 2025-2026 Sở GD&ĐT TP.HCM – Quận 12 - Đề 1 (Bài 2)
Cho phương trình bậc hai $3x^2 - 5x - 4 = 0$.
\begin{enumerate}[a)]
    \item Chứng minh phương trình luôn có hai nghiệm phân biệt $x_1, x_2$.
    \item Không giải phương trình, hãy tính giá trị của biểu thức $P = x_1^2 + x_2^2 + x_1 x_2(2 + 2024x_1 + 2024x_2)$.
\end{enumerate}
\loigiai{
\begin{enumerate}[a)]
    \item Phương trình $3x^2 - 5x - 4 = 0$ có các hệ số $a = 3$, $c = -4$.
    Vì $a \cdot c = 3 \cdot (-4) = -12 < 0$ nên phương trình luôn có hai nghiệm phân biệt trái dấu $x_1, x_2$.
    \item Áp dụng hệ thức Viète cho phương trình $3x^2 - 5x - 4 = 0$, ta có:
    $\begin{cases} S = x_1 + x_2 = -\dfrac{b}{a} = \dfrac{5}{3} \\ P_0 = x_1 x_2 = \dfrac{c}{a} = -\dfrac{4}{3}. \end{cases}$
    Ta có:
    $x_1^2 + x_2^2 = (x_1 + x_2)^2 - 2x_1 x_2 = S^2 - 2P_0 = \left(\dfrac{5}{3}\right)^2 - 2 \cdot \left(-\dfrac{4}{3}\right) = \dfrac{25}{9} + \dfrac{8}{3} = \dfrac{49}{9}$.
    Mặt khác:
    $x_1 x_2(2 + 2024x_1 + 2024x_2) = x_1 x_2 \left[2 + 2024(x_1 + x_2)\right] = P_0(2 + 2024S)$.
    Thay số:
    $P_0(2 + 2024S) = -\dfrac{4}{3} \cdot \left(2 + 2024 \cdot \dfrac{5}{3}\right) = -\dfrac{4}{3} \cdot \left(\dfrac{6 + 10120}{3}\right) = -\dfrac{4}{3} \cdot \dfrac{10126}{3} = -\dfrac{40504}{9}$.
    Do đó giá trị biểu thức $P$ là:
    $P = \dfrac{49}{9} + \left(-\dfrac{40504}{9}\right) = -\dfrac{40455}{9} = -4495$.
\end{enumerate}
}
\end{bt}

\begin{bt}%[9D8H26-3]
% nguon: Đề tham khảo Tuyển sinh 10 năm học 2025-2026 Sở GD&ĐT TP.HCM – Quận 12 - Đề 1 (Bài 3)
Biểu đồ theo dõi lượng mưa trung bình các tháng trong năm tại Hà Nội và Thành phố Hồ Chí Minh cho thấy: Thành phố Hồ Chí Minh có mùa mưa kéo dài trong $6$ tháng (từ tháng 5 đến tháng 10) với lượng mưa hàng tháng đều vượt mức $200\text{ mm}$, trong khi $6$ tháng còn lại có lượng mưa dưới $200\text{ mm}$. Chọn ngẫu nhiên một tháng trong năm. Biết rằng khả năng chọn mỗi tháng là như nhau.
\begin{enumerate}[a)]
    \item Tính xác suất của biến cố $P$: ``Tháng được chọn có lượng mưa trung bình của Thành phố Hồ Chí Minh lớn hơn $200\text{ mm}$''.
    \item Biết rằng trong $12$ tháng của năm, có đúng $5$ tháng mà độ chênh lệch lượng mưa trung bình giữa Hà Nội và Thành phố Hồ Chí Minh không vượt quá $40\text{ mm}$. Tính xác suất của biến cố $Q$: ``Tháng được chọn có lượng mưa chênh lệch giữa hai thành phố không vượt quá $40\text{ mm}$''.
\end{enumerate}
\loigiai{
Một năm có $12$ tháng, do đó số kết quả có thể xảy ra khi chọn ngẫu nhiên một tháng là $n(\Omega) = 12$.
\begin{enumerate}[a)]
    \item Có $6$ tháng (từ tháng 5 đến tháng 10) mà lượng mưa trung bình của Thành phố Hồ Chí Minh lớn hơn $200\text{ mm}$, do đó số kết quả thuận lợi cho biến cố $P$ là $n(P) = 6$.
    Xác suất của biến cố $P$ là:
    $P(P) = \dfrac{n(P)}{n(\Omega)} = \dfrac{6}{12} = \dfrac{1}{2} = 0{,}5$.
    \item Có $5$ tháng mà lượng mưa chênh lệch giữa Hà Nội và Thành phố Hồ Chí Minh không vượt quá $40\text{ mm}$, do đó số kết quả thuận lợi cho biến cố $Q$ là $n(Q) = 5$.
    Xác suất của biến cố $Q$ là:
    $P(Q) = \dfrac{n(Q)}{n(\Omega)} = \dfrac{5}{12}$.
\end{enumerate}
}
\end{bt}

\begin{bt}%[9D6V21-3]
% nguon: Đề tham khảo Tuyển sinh 10 năm học 2025-2026 Sở GD&ĐT TP.HCM – Quận 12 - Đề 1 (Bài 4)
Một mảnh vườn hình chữ nhật có chiều rộng bằng $a$ (m) và chiều dài hơn chiều rộng $6\text{ m}$. Người ta bớt ra hai dải đất chạy dọc theo mép mảnh vườn có bề rộng là $x$ (m) để làm lối đi (với $0 < x < a$).
\begin{enumerate}[a)]
    \item Viết biểu thức biểu thị diện tích phần đất còn lại của mảnh vườn để trồng trọt theo $a$ và $x$.
    \item Tính diện tích phần đất còn lại khi $a = 30\text{ m}$ và $x = 1\text{ m}$.
\end{enumerate}
\loigiai{
\begin{enumerate}[a)]
    \item Chiều rộng ban đầu của mảnh vườn là $a$ (m), chiều dài ban đầu là $a + 6$ (m).
    Sau khi bớt lối đi rộng $x$ (m), kích thước phần đất còn lại là:
    Chiều rộng còn lại: $a - x$ (m);
    Chiều dài còn lại: $(a + 6) - x$ (m).
    Diện tích phần đất còn lại của mảnh vườn là:
    $S = (a - x)(a + 6 - x)\text{ (m}^2\text{)}$.
    \item Khi $a = 30\text{ m}$ và $x = 1\text{ m}$, ta có:
    Chiều rộng còn lại: $30 - 1 = 29\text{ m}$.
    Chiều dài còn lại: $(30 + 6) - 1 = 35\text{ m}$.
    Diện tích phần đất còn lại là:
    $S = 29 \cdot 35 = 1015\text{ m}^2$.
\end{enumerate}
}
\end{bt}

\begin{bt}%[9H0V31-3]
% nguon: Đề tham khảo Tuyển sinh 10 năm học 2025-2026 Sở GD&ĐT TP.HCM – Quận 12 - Đề 1 (Bài 5)
Một chiếc bình chứa nước có lòng trong dạng hình trụ với chiều cao $20\text{ cm}$ và chu vi đường tròn đáy là $10\pi\text{ cm}$.
\begin{enumerate}[a)]
    \item Tính thể tích nước tối đa chiếc bình có thể chứa khi đổ đầy (kết quả làm tròn đến hai chữ số thập phân của $\text{cm}^3$).
    \item Hiện tại lượng nước trong bình đang có mực nước cao $10\text{ cm}$. Một con quạ muốn uống nước trong bình nên đã gắp những viên sỏi có dạng hình cầu đường kính $4\text{ cm}$ thả vào bình. Biết mỏ con quạ chỉ có thể với tới khoảng cách tối đa $6\text{ cm}$ kể từ miệng bình. Hỏi con quạ cần thả tối thiểu bao nhiêu viên sỏi như vậy vào bình để có thể uống được nước?
\end{enumerate}
\loigiai{
\begin{enumerate}[a)]
    \item Bán kính đáy của lòng bình hình trụ là:
    $R = \dfrac{C}{2\pi} = \dfrac{10\pi}{2\pi} = 5\text{ cm}$.
    Thể tích nước tối đa bình có thể chứa là:
    $V = \pi R^2 h = \pi \cdot 5^2 \cdot 20 = 500\pi \approx 1570{,}80\text{ cm}^3$.
    \item Chiều cao của bình là $20\text{ cm}$, mỏ quạ với tới tối đa cách miệng bình $6\text{ cm}$, do đó mực nước trong bình cần dâng lên đến độ cao tối thiểu là:
    $h_{\min} = 20 - 6 = 14\text{ cm}$.
    Mực nước hiện tại cao $10\text{ cm}$, do đó mực nước cần dâng thêm ít nhất một đoạn là:
    $\Delta h = 14 - 10 = 4\text{ cm}$.
    Thể tích nước cần dâng thêm tối thiểu là:
    $V_{\text{dâng}} = \pi R^2 \Delta h = \pi \cdot 5^2 \cdot 4 = 100\pi\text{ cm}^3$.
    Mỗi viên sỏi hình cầu có đường kính $4\text{ cm}$ nên có bán kính $r = \dfrac{4}{2} = 2\text{ cm}$.
    Thể tích của một viên sỏi là:
    $V_{\text{sỏi}} = \dfrac{4}{3}\pi r^3 = \dfrac{4}{3}\pi \cdot 2^3 = \dfrac{32\pi}{3}\text{ cm}^3$.
    Số viên sỏi con quạ cần thả vào bình tối thiểu là:
    $n = \dfrac{V_{\text{dâng}}}{V_{\text{sỏi}}} = \dfrac{100\pi}{\frac{32\pi}{3}} = \dfrac{300}{32} = 9{,}375$.
    Vì số viên sỏi phải là số nguyên dương nên con quạ cần thả tối thiểu $10$ viên sỏi để có thể uống được nước.
\end{enumerate}
}
\end{bt}

\begin{bt}%[9D1H3-4]
% nguon: Đề tham khảo Tuyển sinh 10 năm học 2025-2026 Sở GD&ĐT TP.HCM – Quận 12 - Đề 1 (Bài 6)
Formalin là dung dịch chứa từ $37\%$ đến $40\%$ formaldehyde, có tác dụng kháng khuẩn và kháng nấm nên thường được dùng làm chất bảo quản y tế. Một nhà máy sản xuất hóa chất đang có sẵn dung dịch formaldehyde nồng độ $15\%$ và dung dịch formaldehyde nồng độ $65\%$.
\begin{enumerate}[a)]
    \item Tính thể tích mỗi loại dung dịch formaldehyde trên cần dùng để pha chế được $300$ lít dung dịch formaldehyde nồng độ $35\%$ (giả sử lượng hóa chất không hao hụt).
    \item Một cơ sở y tế đặt hàng nhà máy một lô dung dịch formalin. Nhà máy sử dụng $200$ lít dung dịch formaldehyde $15\%$ trộn với một lượng dung dịch formaldehyde $65\%$ để sản xuất formalin. Hỏi thể tích dung dịch formaldehyde $65\%$ cần thêm vào nằm trong khoảng nào để sản phẩm đạt tiêu chuẩn formalin (kết quả làm tròn đến hàng đơn vị của lít)?
\end{enumerate}
\loigiai{
\begin{enumerate}[a)]
    \item Gọi thể tích dung dịch formaldehyde $15\%$ và $65\%$ cần dùng lần lượt là $x$ và $y$ (lít), điều kiện $x, y > 0$.
    Vì tổng thể tích dung dịch thu được là $300$ lít nên ta có phương trình:
    $x + y = 300$. \quad (1)
    Lượng formaldehyde nguyên chất có trong $300$ lít dung dịch $35\%$ là:
    $0{,}35 \cdot 300 = 105$ (lít).
    Ta có phương trình:
    $0{,}15x + 0{,}65y = 105$. \quad (2)
    Từ (1) và (2) ta có hệ phương trình:
    $\begin{cases} x + y = 300 \\ 0{,}15x + 0{,}65y = 105 \end{cases} \iff \begin{cases} 0{,}15x + 0{,}15y = 45 \\ 0{,}15x + 0{,}65y = 105 \end{cases} \iff \begin{cases} 0{,}5y = 60 \\ x = 300 - y \end{cases} \iff \begin{cases} y = 120 \\ x = 180 \end{cases}$ (thỏa mãn điều kiện).
    Vậy cần dùng $180$ lít dung dịch formaldehyde $15\%$ và $120$ lít dung dịch formaldehyde $65\%$.
    \item Gọi $V$ là thể tích dung dịch formaldehyde $65\%$ cần thêm vào ($V > 0$, đơn vị: lít).
    Tổng thể tích dung dịch sau khi pha trộn là $200 + V$ (lít).
    Lượng formaldehyde nguyên chất có trong hỗn hợp là:
    $200 \cdot 15\% + V \cdot 65\% = 30 + 0{,}65V$ (lít).
    Nồng độ phần trăm của dung dịch thu được là:
    $C\% = \dfrac{30 + 0{,}65V}{200 + V}$.
    Để sản phẩm thu được đạt tiêu chuẩn formalin ($37\% \le C\% \le 40\%$), ta có:
    $0{,}37 \le \dfrac{30 + 0{,}65V}{200 + V} \le 0{,}40$.
    Vì $200 + V > 0$ nên bất đẳng thức tương đương với hệ:
    $\begin{cases} 30 + 0{,}65V \ge 0{,}37(200 + V) \\ 30 + 0{,}65V \le 0{,}40(200 + V) \end{cases} \iff \begin{cases} 30 + 0{,}65V \ge 74 + 0{,}37V \\ 30 + 0{,}65V \le 80 + 0{,}40V \end{cases} \iff \begin{cases} 0{,}28V \ge 44 \\ 0{,}25V \le 50 \end{cases} \iff \begin{cases} V \ge \dfrac{44}{0{,}28} \approx 157{,}14 \\ V \le \dfrac{50}{0{,}25} = 200. \end{cases}$
    Làm tròn đến hàng đơn vị của lít, thể tích dung dịch formaldehyde $65\%$ cần dùng nằm trong khoảng từ $157$ lít đến $200$ lít.
\end{enumerate}
}
\end{bt}

\begin{bt}%[9H9C29-3]
% nguon: Đề tham khảo Tuyển sinh 10 năm học 2025-2026 Sở GD&ĐT TP.HCM – Quận 12 - Đề 1 (Bài 7)
Cho điểm $M$ nằm bên ngoài đường tròn $(O; R)$. Qua $M$ kẻ hai tiếp tuyến $MA, MB$ với đường tròn $(O)$ ($A, B$ là các tiếp điểm). Kẻ đường kính $BD$ của đường tròn $(O)$, đường thẳng $MD$ cắt đường tròn $(O)$ tại điểm thứ hai là $C$. Gọi $H$ là giao điểm của $MO$ và đoạn thẳng $AB$.
\begin{enumerate}[a)]
    \item Chứng minh tứ giác $MAOB$ nội tiếp và $AB \perp OM$.
    \item Chứng minh $MA^2 = MC \cdot MD$ và $MH \cdot MO = MC \cdot MD$.
    \item Cho biết $AB = R\sqrt{3}$. Tính độ dài đoạn thẳng $MH$ theo $R$.
\end{enumerate}
\loigiai{
\begin{enumerate}[a)]
    \item Vì $MA, MB$ là hai tiếp tuyến của đường tròn $(O)$ nên $MA \perp OA$ và $MB \perp OB$, suy ra $\widehat{MAO} = 90^\circ$ và $\widehat{MBO} = 90^\circ$.
    Tứ giác $MAOB$ có $\widehat{MAO} + \widehat{MBO} = 90^\circ + 90^\circ = 180^\circ$ nên là tứ giác nội tiếp đường tròn đường kính $MO$.
    Lại có $MA = MB$ (tính chất hai tiếp tuyến cắt nhau) và $OA = OB = R$, suy ra $MO$ là đường trung trực của đoạn thẳng $AB$, do đó $AB \perp OM$ tại $H$.
    \item Xét $\triangle MAC$ và $\triangle MDA$, ta có:
    Góc $\widehat{AMD}$ chung;
    $\widehat{MAC} = \widehat{MDA}$ (góc tạo bởi tiếp tuyến và dây cung và góc nội tiếp cùng chắn cung $AC$).
    Do đó $\triangle MAC \backsim \triangle MDA$ (g.g), suy ra $\dfrac{MA}{MD} = \dfrac{MC}{MA} \implies MA^2 = MC \cdot MD$.
    Mặt khác, trong tam giác vuông $MAO$ tại $A$ có đường cao $AH$, theo hệ thức lượng ta có:
    $MA^2 = MH \cdot MO$.
    Từ đó suy ra $MH \cdot MO = MC \cdot MD$ (đpcm).
    \item Vì $MO$ là đường trung trực của $AB$ nên $H$ là trung điểm của $AB$, do đó:
    $AH = \dfrac{AB}{2} = \dfrac{R\sqrt{3}}{2}$.
    Trong tam giác vuông $OAH$ vuông tại $H$ có $OA = R$:
    $OH = \sqrt{OA^2 - AH^2} = \sqrt{R^2 - \left(\dfrac{R\sqrt{3}}{2}\right)^2} = \sqrt{R^2 - \dfrac{3R^2}{4}} = \sqrt{\dfrac{R^2}{4}} = \dfrac{R}{2}$.
    Trong tam giác vuông $MAO$ vuông tại $A$ có đường cao $AH$:
    $AH^2 = MH \cdot OH \implies MH = \dfrac{AH^2}{OH} = \dfrac{\left(\dfrac{R\sqrt{3}}{2}\right)^2}{\dfrac{R}{2}} = \dfrac{\dfrac{3R^2}{4}}{\dfrac{R}{2}} = \dfrac{3R}{2}$.
    Vậy $MH = \dfrac{3R}{2}$.
\end{enumerate}
}
\end{bt}
